\subsection{正规积分子群}

引理 3. 设 \(\mathbf{G}\) 是李群，\(\mathrm{H}_1\) 和 \(\mathrm{H}_2\) 是拓扑有可数基的积分子群，且 \(g \in \mathbf{G}\)。则

\[
g \mathrm{H} _ {1} g ^ {- 1} = \mathrm{H} _ {2} \Leftrightarrow (\operatorname{Ad} g) (\mathrm{L} (\mathrm{H} _ {1})) = \mathrm{L} (\mathrm{H} _ {2}).
\]

\(\operatorname{Ad} g = \mathrm{T}_{e}(\operatorname{Int} g)\)。因此，由结构传递，\((\operatorname{Int} g)(\mathrm{H}_{1})\) 的李代数为 \((\operatorname{Ad} g)(\mathrm{L}(\mathrm{H}_{1}))\)。由于 \(\mathrm{H}_{1}\) 和 \(\mathrm{H}_{2}\) 都有可数基，说集合 \(\mathrm{H}_{2}\) 与 \((\operatorname{Int} g)(\mathrm{H}_{1})\) 相等，等价于说积分子群 \(\mathrm{H}_{2}\) 与 \((\operatorname{Int} g)(\mathrm{H}_{1})\) 相等（no. 2, Corollary 3 to Proposition 3）。于是引理由定理 2 (i) 得出。

命题 14. 设 G 是李群，H 是拓扑有可数基的积分子群。以下条件等价：

\begin{enumerate}
\def\labelenumi{(\roman{enumi})}
\item
  H 是 G 的正规子群；
\item
  L(H) 在 Ad(G) 下不变。
\end{enumerate}

若进一步假设 G 连通，则这些条件还等价于：

\begin{enumerate}
\def\labelenumi{(\roman{enumi})}
\setcounter{enumi}{2}
\tightlist
\item
  L(H) 是 L(G) 的理想。
\end{enumerate}

若进一步假设 \(\mathbf{G}\) 单连通且 \(\mathbf{L}(\mathbf{H})\) 在 \(\mathbf{L}(\mathbf{G})\) 中具有有限余维，则这些条件蕴含 \(\mathbf{H}\) 是 \(\mathbf{G}\) 的李子群，且 \(\mathbf{G} / \mathbf{H}\) 单连通。

(i) \(\Leftrightarrow\) (ii) 的等价性由引理 3 得出。若进一步假设 G 连通，则条件 (ii) 等价于 L(H) 在 ad L(G) 下稳定（no. 5, Corollary 1 to Proposition 13 and § 3, no. 12, Proposition 44）。

假设 G 单连通且 L(H) 是 L(G) 中有限余维的理想。由第 3 号的 Theorem 3，存在单连通李群 G'，使 L(G') 同构于 L(G)/L(H)。存在一个从 L(G) 到 L(G') 的连续满态射 f，核为 L(H)。由第 1 号的 Theorem 1，存在从 G 到 G' 的态射 \(\phi\)，使 L( \(\phi\) ) = f.~该态射是下沉映射，因而其核 N 是 G 的李子群，且 L(N) = Ker f = L(H)。所以 H 是 N 的单位元分支，因而是 G 的李子群。令 \(\psi\) 为将 \(\phi\) 通过商得到的从 G/H 到 G' 的态射。该态射是 étale 的；由于 G 单连通，\(\psi\) 是从 G/H 到 G' 的同构。

推论 1. 设 \(G\) 是有限维单连通李群。令 \(m, h\) 为 \(L(G)\) 的李子代数，使得 \(L(G)\) 是 \(m\) 与 \(h\) 的半直积。令 \(M, H\) 为 \(G\) 中相应的积分子群。则 \(M\) 和 \(H\) 是 \(G\) 的单连通李子群，并且作为李群，\(G\) 是 \(M\) 与 \(H\) 的半直积。

由命题 14，H 是 G 的正规李子群，且李群 G/H 单连通。令 \(\pi\) 为从 G 到 G/H 的规范态射。存在从 G/H 到 M 的态射 \(\theta\)，使 L(θ) 是从 L(G)/L(H) = L(G)/h 到 L(M) = m 的规范同构。于是

\[
\mathrm{L} (\pi \circ \theta) = \mathrm{L} (\pi) \circ \mathrm{L} (\theta) = \mathrm{Id} _ {\mathrm{L} (\mathrm{G} / \mathrm{H})}
\]

从而 \(\pi \circ \theta = \mathrm{Id}_{\mathrm{G}/\mathrm{H}}\)。由第 1 号 Proposition 1 的 Corollary 1，\(\theta (\mathrm{G / H}) = \mathbf{M}\)。由 § 1, no. 4 的 Proposition 8，M 是 G 的李子群，且李群 G 是 M 与 H 的半直积。

推论 2. 设 G 是有限维单连通李群，H 是 G 的正规连通李子群，\(\pi\) 是从 G 到 G/H 的规范态射。

\begin{enumerate}
\def\labelenumi{(\roman{enumi})}
\item
  存在从 G/H 到 G 的解析映射 \(\rho\)，使得 \(\pi \circ \rho = \mathrm{Id}_{\mathrm{G / H}}\)。
\item
  对每个具有 (i) 中性质的映射 \(\rho\)，映射 \((h, m) \mapsto h\rho(m)\) 从 \(\mathbf{H} \times (\mathbf{G}/\mathbf{H})\) 到 \(\mathbf{G}\)，是解析流形同构。
\item
  H 和 G/H 都是单连通的。
\end{enumerate}

令 \(n = \dim G - \dim H\)。当 \(n = 0\) 时推论显然成立。我们对 \(n\) 作归纳论证。

假设存在一个包含 L(H) 且不同于 L(G) 和 L(H) 的 L(G) 理想。令 H' 为相应的 G 的连通李子群。令 \(\pi_{1}: G \to G/H'\) 和 \(\pi_{2}: H' \to H'/H\) 为规范态射。由归纳假设，存在解析映射 \(\rho_{1}: G/H' \to G\)、\(\rho_{2}: H'/H \to H'\)，使得 \(\pi_{1} \circ \rho_{1} = Id_{G/H'}\)、\(\pi_{2} \circ \rho_{2} = Id_{H'/H}\)。令

\[
\pi_ {3} \colon \mathrm{G} / \mathrm{H} \rightarrow \mathrm{G} / \mathrm{H} ^ {\prime}
\]

为规范态射。若 \(x \in \mathrm{G} / \mathrm{H}\)，且 \(y\) 是 \(x\) 在 \(\mathrm{G}, y\) 中的代表元，则 y 和 \(\rho_1(\pi_3(x))\) 模 \(\mathrm{H}'\) 有相同的规范像，因而 \(x^{-1}\pi (\rho_1(\pi_3(x))) \in \mathrm{H}' / \mathrm{H}\)。令

\[
\rho (x) = \rho_ {1} (\pi_ {3} (x)) \rho_ {2} (\pi (\rho_ {1} (\pi_ {3} (x))) ^ {- 1} x) \in G.
\]

显然，\(\rho\) 是从 G/H 到 G 的解析映射，并且

\[
\begin{array}{r l} \pi (\rho (x)) & = \pi (\rho_ {1} (\pi_ {3} (x))) \pi_ {2} (\rho_ {2} (\pi (\rho_ {1} (\pi_ {3} (x))) ^ {- 1} x)) \\ & = \pi (\rho_ {1} (\pi_ {3} (x))) \pi (\rho_ {1} (\pi_ {3} (x))) ^ {- 1} x = x. \end{array}
\]

如果现在 \(\mathbf{L}(\mathbf{G})\) 的理想中，包含 \(\mathbf{L}(\mathbf{H})\) 的只有 \(\mathbf{L}(\mathbf{G})\) 和 \(\mathbf{L}(\mathbf{H})\)，则李代数 \(\mathrm{L(G) / L(H)}\) 要么是一维的，要么是单的。在这两种情形下，\(\mathrm{L(G)}\) 都是某个子代数与 \(\mathrm{L(H)}\) 的半直积；第一种情形显然，第二种情形由 Chapter I, § 6, Corollary 3 to Theorem 5 得出。因此断言 (i) 由推论 1 得出。

断言 (ii) 显然。断言 (iii) 由 (i) 和 (ii) 得出。

当 G 是无限维的或 H 不是正规子群时，推论 2 的结论不再必然成立（练习 8 和 15）。

推论 3. 设 \(\mathbf{G}\) 是有限维连通李群，\(\mathbf{H}\) 是 \(\mathbf{G}\) 的正规连通李子群。从 \(\pi_1(\mathbf{H})\) 到 \(\pi_1(\mathbf{G})\) 的规范态射是单射。

令 \(G_{1}\) 为 G 的万有覆盖，\(\lambda\) 为从 \(G_{1}\) 到 G 的规范映射。将 \(G_{1}\) 的李代数与 L(G) 识别。\(\lambda^{-1}(H)\) 是 \(G_{1}\) 的李子群，在 \(G_{1}\) 中正规，且其李代数为 L(H)。令 \(H_{1}\) 为 \(\lambda^{-1}(H)\) 的单位元分支，\(\lambda_{1} = \lambda | H_{1}\)。则 \(\mathrm{L}(H_{1}) = \mathrm{L}(H)\)，从而 \(\lambda_{1}\) 是从 \(H_{1}\) 到 H 的 étale 态射。另一方面，\(H_{1}\) 是单连通的（推论 2），因而可与 H 的万有覆盖识别。于是根据 General Topology, Chapter XI，从 \(\pi_{1}(H)\) 到 \(\pi_{1}(G)\) 的规范态射可识别为从 Ker \(\lambda_{1}\) 到 Ker \(\lambda\) 的规范嵌入。

